\chapter{Método dos Elementos Finitos}

% Configurações de documento e pacotes equivalentes




\begin{figure}[htbp]
  \centering
  \begin{tikzpicture}[
      node distance=1.5cm and 1.0cm,
      block/.style={rectangle, draw=blue!70, fill=blue!5, rounded corners=5pt, thick, align=center, inner sep=6pt},
      line/.style={draw, -latex, thick, color=gray!80}
    ]
    \node [block, width=6cm] (step1) {\textbf{1. Definição do Problema} \\ \small Vaso de pressão multicamadas};
    \node [block, width=5.5cm, below=0.8cm of step1] (step2) {\textbf{2. Parâmetros Iniciais} \\ \small Geometria, Fração volumétrica, Condições de contorno};
    \node [block, width=4.8cm, below left=1.0cm and 0.2cm of step2] (step3a) {\textbf{3a. Modelo Analítico} \\ \small Homogeneização Assintótica (MHA)};
    \node [block, width=4.8cm, below right=1.0cm and 0.2cm of step2] (step3b) {\textbf{3b. Modelo Numérico} \\ \small Elementos Finitos (MEF -- FEniCS)};
    \node [block, width=5.5cm, below=3.2cm of step2] (step4) {\textbf{4. Verificação} \\ \small Comparação (Tensões e Deslocamentos)};
    \node [block, width=6cm, below=0.8cm of step4] (step5) {\textbf{5. Análise de Viabilidade} \\ \small Resistência ($W_g$) e Rigidez ($K_g$)};

    \draw [line] (step1) -- (step2);
    \draw [line] (step2) -| (step3a);
    \draw [line] (step2) -| (step3b);
    \draw [line] (step3a) |- (step4);
    \draw [line] (step3b) |- (step4);
    \draw [line] (step4) -- (step5);
  \end{tikzpicture}
  \caption{Fluxograma geral da metodologia adotada. Fonte: Autor.}
  \label{fig-fluxograma}
\end{figure}

A modelagem numérica de estruturas elásticas microperiódicas com alto contraste de propriedades exige uma transição rigorosa entre as formulações matemáticas contínuas desenvolvidas no Capítulo 3 e a sua representação no domínio discreto. Este capítulo expõe a metodologia prática e o fluxo computacional adotados para a simulação de vasos de pressão cilíndricos multicamadas por meio do Método dos Elementos Finitos (MEF). São detalhados a estratégia de discretização espacial do domínio axissimétrico, a lógica de programação para a geração da microestrutura bimetálica periódica via FEniCS, o equacionamento das condições de contorno de pressão e simetria longitudinal, e o protocolo de cálculo de métricas de eficiência estrutural baseados no código computacional real de projeto implementado no arquivo principal `MHA vs FEM - main - EDITAR APENAS ESSE.py`.

\section{Discretização e Configuração da Malha}\label{sec:discretizacao_malha}

A precisão de um solucionador de elementos finitos para problemas com heterogeneidades de alta frequência e gradientes de tensão severos é intrinsecamente governada pelas características geométricas do domínio computacional e pelo refinamento espacial da malha de discretização [35]. Esta seção detalha a definição do domínio, a escolha da classe de elementos finitos e o equacionamento para o posicionamento exato dos nós nas interfaces de transição elástica de acordo com a rotina algorítmica real.

\subsection{Geometria do Domínio Computacional Axissimétrico}

Conforme discutido na formulação teórica da elasticidade de revolução (Seção 3.4), a invariância circunferencial ($\partial_{\theta} = 0$) e a hipótese de deformação plana generalizada reduzem o problema elástico tridimensional (3D) de um cilindro oco para uma seção transversal bidimensional (2D) contida no plano radial-axial $r-z$ [190].

Para manter total consistência físico-matemática e dimensional com o código numérico desenvolvido, o domínio computacional espacial, denotado por $Omega_{\text{rz}}$, é representado geometricamente por uma seção retangular plana cujas dimensões de projeto em escala real de simulação são dadas por:

$ \Omega\textit{("rz") = { (r, z) \in \mathbb{R}^2 : R}("int") \le  r \le  R\textit{("ext") , \quad 0 \le  z \le  H } $ <eq:dominio}mef>

onde:
\begin{itemize}
  \item $R_{\text{int}} = 9,0 "m"$ representa o raio interno do cilindro cilíndrico de referência;
  \item $R_{\text{ext}} = 11,0 "m"$ representa o raio externo do cilindro;
  \item A largura do domínio equivale à espessura total da parede do vaso, dada por $t = R\textit{("ext") - R}("int") = 2,0 "m"$;
  \item A dimensão axial (altura do retângulo de revolução) é fixada em $H = 2,0 "m"$.
\end{itemize}

Essa escala geométrica em metros é rigorosamente compatível com as unidades das constantes constitutivas elásticas eloficiais (Pascals, $10^9 "Pa"$) e da pressão nominal aplicada ($30 "MPa"$, ou $3,0 \times 10^7 "Pa"$) utilizadas no script principal, evitando inconsistências de fator de escala numérico durante a montagem das equações elásticas.

\subsection{Escolha do Elemento Finito (Lagrange de Grau 1)}

Na implementação do solucionador no FEniCS, as componentes do vetor de deslocamento $\mathbf{u} = (u\textit{r, u}z)$ e o campo de funções de teste virtuais $\mathbf{v} = (v\textit{r, v}z)$ são interpolados espacialmente utilizando \textbf{elementos finitos de Lagrange de grau 1} ($P_1$), configurando um espaço de funções contínuas de Galerkin (`VectorFunctionSpace(mesh, "CG", 1)`) [926, 1249].

A adoção de polinômios lineares, em conjunto com uma malha altamente refinada e estruturada, é uma premissa metodológica fundamental para a simulação elástica eficiente de cilindros concêntricos sob pressão pelas seguintes razões elásticas e computacionais [1249]:

1. \textbf{Invariância de Deformação e Tensão por Elemento:} Como a deformação $\mathbf{\epsilon}$ depende das derivadas primeiras do campo de deslocamento ($\mathbf{\epsilon} = 1/2(\mathbf{nabla} u + \mathbf{nabla}^T u)$), o uso de funções de interpolação lineares ($P_1$) garante que as deformações e as tensões no interior de cada elemento finito individual sejam de classe $C^{-1}$, ou seja, constantes ao longo de cada triângulo [605]. Isso é perfeitamente congruente com a física de meios laminados periódicos (onde o material e as propriedades elásticas também são assumidos como constantes por trecho/elemento), fornecendo degraus de tensões constantes por subcamada elástica sem introduzir suavizações ou gradientes parabólicos espúrios na interface de descontinuidade elástica [1249].
2. \textbf{Minimização Drástica do Custo Computacional:} Como o presente estudo conduz simulações com um elevado número de células periódicas (até $N = 200$) e resoluções de malha extremamente refinadas, a adoção de elementos $P_1$ limita significativamente o número total de graus de liberdade globais a serem resolvidos. Isso reduz drasticamente o tempo de CPU e os requisitos de memória RAM, permitindo a execução de varreduras paramétricas e tabelas de erros quase que instantaneamente.
3. \textbf{Resolução Radial por Refinamento Dinâmico:} A potencial limitação de menor precisão local de elementos lineares em gradientes acentuados é completamente mitigada pelo emprego de uma altíssima densidade radial de elementos finitos, garantindo que as curvas globais hiperbólicas de deslocamento e tensão sejam perfeitamente discretizadas de forma facetada e suave.

\subsection{Configuração de Alta Discretização e Densidade Radial}

A malha do domínio $\Omega\textit{("rz")$ é estruturada de forma mapeada, contendo $N}z$ divisões na direção axial e $N_r$ divisões na direção radial, resultando em uma malha de elementos triangulares obtidos pela divisão diagonal do retângulo espacial [926, 927]. Essa discretização ao longo do raio $r$ é governada pela resolução espacial requerida para capturar os degraus microscópicos locais de descontinuidade constitutiva nas subcamadas [927].

No algoritmo implementado no arquivo principal `MHA vs FEM - main.py`, a densidade de elementos radiais respeita de forma exata e proporcional as frações volumétricas elásticas de cada camada, definindo-se no script [927]:
\begin{itemize}
  \item `V1`: Fração volumétrica do Material 1;
\end{itemize}

\begin{itemize}
  \item `V2`: Fração volumétrica do Material 2;
\end{itemize}

\begin{itemize}
  \item `N_cel`: Número de células periódicas;
\end{itemize}

\begin{itemize}
  \item `resolucao_malha`: Número de elementos radiais por célula periódica;
\end{itemize}

\begin{itemize}
  \item `nel\textit{1 = int(V1 * resolucao}malha)`: número de elementos na camada do Material 1 (Aço) por célula periódica [927];
\end{itemize}

\begin{itemize}
  \item `nel\textit{2 = int(V2 * resolucao}malha)`: número de elementos na camada do Material 2 (Alumínio) por célula periódica, com $V\textit{2 = 1 - V}1$ [927];
\end{itemize}

\begin{itemize}
  \item `nel\textit{celula = nel}1 + nel_2`: número total de elementos radiais contidos em cada período elástico básico (célula unitária) [927].
\end{itemize}

Dessa forma, o número total de elementos na direção radial ($N\textit{r$, ou `n}elements\textit{r`) e axial ($N}z$, ou `n\textit{elem}z`) são estabelecidos diretamente no gerador `RectangleMesh` por [927]:

$ N\textit{r = N}("cel") \times "resolucao\textit{malha" $ <eq:nr}mef_v3>
$ N\textit{z = n}("elem\textit{z") $ <eq:nz}mef_v3>

```python
# Trecho de discretização espacial do domínio de parede do vaso
nel\textit{1 = int(V1 * resolucao}malha)
nel\textit{2 = int(V2 * resolucao}malha)
nel\textit{celula = nel}1 + nel_2
n\textit{elements}r = N\textit{celulas * nel}celula

mesh = RectangleMesh(
    Point(R\textit{int, 0.0), Point(R}ext, H),
    int(n\textit{elements}r), int(n\textit{elem}z)
)
```

Adotando-se a configuração padrão de simulação com $N\textit{{"cel"} = 200$ células periódicas, `resolucao}malha` = 20 e $N\textit{z = 10$, a malha radial totaliza \textbf{$N}r = 4.000$ elementos radiais}, gerando uma discretização de $80.000$ elementos triangulares no domínio bidimensional [266, 927]. Essa altíssima densidade radial garante que cada subcamada individual de aço ou alumínio (cuja espessura de período para $N=200$ é de apenas $5,0{ "mm"}$ no aço e $5,0{ "mm"}$ no alumínio) contenha exatamente \textbf{10 elementos lineares internos}, permitindo mapear com riqueza de detalhes elásticos as flutuações e tensões microscópicas [63, 831].

\subsection{O Alinhamento dos Nós com as Interfaces Físicas}

No Método dos Elementos Finitos, a distribuição das propriedades mecânicas do material (módulo de elasticidade $E$ e coeficiente de Poisson $\nu$) é avaliada e integrada numericamente nos pontos de integração Gaussiana localizados no interior de cada elemento finito [59].

Se um elemento finito cruzar espacialmente a fronteira física que divide uma camada de aço de uma camada de alumínio, as propriedades elásticas integradas para aquele elemento sofrerão uma média suave artificial ("borramento" interfacial) [59, 60]. Como consequência física direta, o salto descontínuo abrupto na tensão circunferencial ($\sigma\textit{(\theta \theta)$) e a mudança de inclinação (\textbf{kinks}) no deslocamento radial ($u}r$), característicos da interface de contato perfeito bimetálica, seriam suavizados, arruinando completamente a acurácia física e a consistência local da simulação numérica [60, 142].

Para blindar o modelo contra essa perda de precisão, a fronteira de cada subcamada material \textbf{deve coincidir estritamente com as bordas dos elementos finitos} [60]. Ou seja, os nós da malha espacial devem estar localizados exatamente sobre as coordenadas radiais das interfaces elásticas [60, 142].

Essa restrição é perfeitamente satisfeita no código real ao impor o arredondamento inteiro de `nel\textit{1` e `nel}2` por célula periódica através da função `int()` em Python [927]. Ao garantir que $n\textit{{"el}cam1"}, n\textit{{"el}cam2"} \in bb(N)$, a espessura de cada camada física se torna um múltiplo inteiro exato do tamanho local do elemento ($h\textit{r = t / N}r$) [63, 64]. O gerador de malhas realiza a transição de propriedades de forma abrupta a cada nó interfacial, erradicando por completo oscilações numéricas parasitárias próximas às interfaces (fenômeno de Gibbs elástico) e garantindo a convergência estrita da simulação heterogênea para o limite assintótico homogeneizado do MHA [64, 154].

A simulação numérica de meios altamente heterogêneos e microperiódicos, como o cilindro bimetálico composto por camadas alternadas de aço carbono e liga de alumínio, impõe desafios significativos de modelagem geométrica. À medida que o número de períodos $N_{"cell"}$ aumenta, a espessura de cada subcamada individual encolhe de forma inversamente proporcional, tornando a reconstrução física de cada interface no ambiente de CAD tradicional e o posterior corte da malha (\textbf{mesh partitioning}) processos manuais altamente ineficientes, propensos a erros de precisão numérica e impossíveis de serem automatizados de forma paramétrica [47].

Para contornar essa barreira operacional, adota-se neste trabalho a estratégia de \textbf{Mapeamento Implícito de Propriedades} via vetor de interfaces discretas. Em vez de discretizar as interfaces físicas na própria malha, o domínio espacial é particionado por uma malha estruturada única e contínua no plano $r z$ (conforme discutido na \marginpar{\tiny Ajustar seção discretizacao\textit{malha}% \ref{sec:discretizacao}malha}). As propriedades elásticas locais dos materiais de fase (módulo de Young, coeficiente de Poisson e a respectiva matriz constitutiva) são injetadas diretamente nos pontos de integração Gaussiana através de expressões analíticas espaciais periódicas de busca indexada programadas no ambiente FEniCS [562].

\subsection{Formulação Matemática do Mapeamento Periódico}

Considere o cilindro de parede espessa definido no intervalo radial $r \in [R\textit{{"int"}, R}{"ext"}]$, cuja espessura total de parede é $t = R\textit{{"ext"} - R}{"int"}$. O cilindro é composto por $N_{"cell"}$ células periódicas unitárias (RVEs) [562].

A espessura de uma célula periódica unitária (período $Y$) é calculada a partir do número de células e da espessura total do vaso por [562]:

$ L\textit{c = frac(R}{"ext"} - R\textit{{"int"}, N}{"cell"}) $ <eq:L_c>

Sendo $V\textit{1$ e $V}2 = 1 - V_1$ as frações volumétricas dos Materiais 1 (Aço) e 2 (Alumínio), respectivamente, as espessuras individuais de cada subcamada dentro de uma célula são expressas por [562]:

$ h\textit{1 = V}1 L\textit{c , \quad h}2 = V\textit{2 L}c $ <eq:h_camadas>

A partição espacial radial do cilindro é dividida de forma discreta em um vetor de raios interfaciais ${r\textit{k}$, pré-calculado e alinhado de forma exata com as coordenadas de nós correspondentes da malha espacial. Para cada célula periódica $i \in {0, 1, \dots, N}{"cel"} - 1}$, as posições das interfaces são calculadas algoritmicamente por [562]:

\begin{itemize}
  \item Interface interna da célula $i$ (Material 1 para Material 2):
  $ r\textit{{(2i+1)} = R}{"int"} + i L\textit{c + h}1 $ <eq:interface_interna>
\end{itemize}

\begin{itemize}
  \item Interface de contorno entre a célula $i$ e a célula $i+1$ (Material 2 para o Material 1 da célula subsequente):
  $ r\textit{{(2i+2)} = R}{"int"} + (i + 1) L\textit{c $ <eq:interface}externa>
\end{itemize}

onde o índice de célula $i$ estende-se até $N\textit{{"cel"} - 2$ para a interface de transição de células, uma vez que a última face externa coincide com o limite do raio externo $R}{"ext"}$ do cilindro [562]. Essa formulação genérica permite simular de forma robusta e rigorosa frações volumétricas quaisquer das fases, sem limitações de simetria geométrica elástica [562].

A função degrau discreta que governa a distribuição espacial do Módulo de Young local $E(r)$ e do Coeficiente de Poisson local $\nu(r)$ ao longo da parede do cilindro baseia-se na identificação de qual intervalo de subcamada indexada $k$ contém a coordenada do ponto de integração Gaussiana $r = x[0]$ [25, 562]:

\begin{equation}
  E(r) = cases(
    E_1 "se" k "é par (Aço)",
    E_2 "se" k "é ímpar (Alumínio)"
  )
  \label{eq:mapeamento\textit{E}v2}
\end{equation}

\begin{equation}
  \nu(r) = cases(
    nu_1 "se" k "é par (Aço)",
    nu_2 "se" k "é ímpar (Alumínio)"
  )
  \label{eq:mapeamento\textit{nu}v2}
\end{equation}

\subsection{Algoritmo de Implementação via `UserExpression` no FEniCS}

No ambiente computacional FEniCS, o mapeamento espacial contínuo expresso pelas \eqref{eq:mapeamento\textit{E}v2} e \eqref{eq:mapeamento\textit{nu}v2} é implementado de forma nativa e extremamente eficiente através de subclasses da classe base `UserExpression` [562]. Esta abordagem baseada em pesquisa indexada permite que as propriedades elásticas sejam avaliadas em tempo real durante o processo de montagem da matriz de rigidez global, eliminando gargalos de memória [35, 562].

A classe customizada desenvolvida para este propósito, denominada `PeriodicProperty`, é estruturada em Python conforme a rotina real de cálculo implementada no script [25, 562]:

```python
class PeriodicProperty(UserExpression):
    def _\textit{init}_(self, val1, val2, interfaces, **kwargs):
        super()._\textit{init}_(**kwargs)
        self.val1, self.val2 = val1, val2 # Propriedades elásticas das fases
        self.interfaces = interfaces      # Vetor de coordenadas das interfaces radiais

    def eval(self, values, x):
        idx = 0                           # Inicializa o índice da camada como par
        # Varredura sequencial para localizar a camada correspondente
        for r_int in self.interfaces:
            if x[0] <= r_int:
                break
            idx += 1

        # Atribuição condicional baseada na paridade do índice do segmento
        values[0] = self.val1 if idx % 2 == 0 else self.val2

    def value_shape(self):
        return () # Retorna um escalar (campo de propriedades elásticas)
```

Durante a execução do solver, as instâncias de `PeriodicProperty` são utilizadas para construir os campos escalares do Módulo de Young e do Coeficiente de Poisson no domínio discreto, escrevendo-se [562]:

```python
E_func = PeriodicProperty(E1, E2, interfaces, degree=0)
nu_func = PeriodicProperty(nu1, nu2, interfaces, degree=0)
```

Adota-se o parâmetro `degree=0` (interpolação constante por trecho) para garantir que a propriedade mude de forma abrupta e descontínua exatamente nas coordenadas das interfaces físicas, simulando com fidelidade a transição bimetálica real sem suavizações artificiais indesejadas [35, 562].

\subsection{Transição Sem Costura entre o Modelo Detalhado e o Modelo Homogeneizado (MHA)}

A grande elegância metodológica desta arquitetura de programação é a facilidade de transição entre o modelo microscópico detalhado (altamente heterogêneo e oscilatório) e o modelo macroscópico equivalente homogeneizado pelo MHA [47, 562].

No FEniCS, as formas bilinear $a(u,v)$ e linear $L(v)$ são acopladas ao tensor constitutivo de quarta ordem $cal{C}$ [35, 562]. Para alternar as análises mecânicas, o operador computacional simplesmente substitui o argumento de rigidez elástica [562]:

1. \textbf{Modelo Detalhado (MEF-D):} O tensor constitutivo local $cal{C}(r)$ é calculado em cada ponto de integração a partir dos campos oscilatórios $E(r)$ e $\nu(r)$ fornecidos pela classe `PeriodicProperty`, gerando o comportamento real em zigue-zague dos deslocamentos e as oscilações de tensão [47, 562].
2. \textbf{Modelo Homogeneizado (MHA):} O tensor constitutivo é inicializado diretamente como uma matriz ortotrópica constante $hat{cal{C}}$, cujos coeficientes efetivos ($hat{C}\textit{{11}$, $hat{C}}{12}$, $hat{C}_{22}$ e $hat{G}$) são calculados analiticamente através das equações de Backus deduzidas no Capítulo 3 e aplicadas no script [47, 562].

Esta simplicidade operacional e algorítmica demonstra o elevado nível de maturidade computacional da simulação, reduzindo o tempo de preparação do modelo e viabilizando estudos paramétricos de otimização estrutural de vasos de pressão multicamadas [47, 562].

\section{Propriedades dos Materiais e Cenários de Carga}\label{sec:propriedades_cenarios}

Para a execução das simulações numéricas via Método dos Elementos Finitos (MEF) e a consequente avaliação da precisão do Método de Homogeneização Assintótica (MHA), torna-se indispensável definir com precisão os parâmetros físicos e mecânicos que caracterizam as fases constituintes do cilindro bimetálico. Esta seção detalha as propriedades elásticas, as densidades volumétricas, os limites de escoamento adotados para o aço e o alumínio, bem como justifica a escolha do cenário de carregamento termomecânico imposto ao modelo [47, 48].

\subsection{Propriedades Mecânicas e Físicas das Fases Constituintes}

O cilindro periódico analisado é composto pelo par metálico bimetálico formado por aço carbono e liga de alumínio, ambos tratados individualmente como meios homogêneos e isotrópicos em nível de subcamada microestrutural [47]. As propriedades mecânicas de entrada (módulo de Young $E$ e coeficiente de Poisson $\nu$), associadas às propriedades físicas de massa específica ($\rho$) e resistência mecânica elástica limite ($\sigma\textit{y$), são sumarizadas na \cite{tab:propriedades}fases} [47, 48]:

#figure(
  caption: [Propriedades elásticas, físicas e limites de resistência das fases constituintes do cilindro bimetálico conforme implementação no script de projeto.],
  table(
    columns: (1.5fr, 1.0fr, 1.2fr, 1.3fr),
    align: (left, center, center, center),
    stroke: 0.5pt + luma(150),
    fill: (x, y) => if y == 0 { rgb("e6edf5") } else if calc.even(y) { rgb("f8fafc") } else { none },
    [\textbf{Propriedade}], [\textbf{Símbolo}], [\textbf{Fase 1: Aço}], [\textbf{Fase 2: Alumínio}],
    [Módulo de Elasticidade], [$E$], [$200 "GPa"$ ($2,0 \times 10^11 "Pa"$)], [$70 "GPa"$ ($7,0 \times 10^10 "Pa"$)],

    [Coeficiente de Poisson], [$\nu$], [$0,30$], [$0,33$],
    [Densidade Volumétrica], [$\rho$], [$7800 "kg/m"^3$], [$2700 "kg/m"^3$],
    [Limite de Escoamento],
    [$sigma_y$],
    [$400,0 "MPa"$ ($4,0 \times 10^8 "Pa"$)],
    [$270,0 "MPa"$ ($2,7 \times 10^8 "Pa"$)],
  ),
) <tab:propriedades_fases> #margin-note([Isso é um quadro, não uma tabela])

A escolha desses materiais reflete um cenário clássico de alto contraste de rigidez elástica (com a razão de módulos $E\textit{{{"aço"}}/E}{{"alu"}} approx 2,86$), o qual representa o cenário ideal e mais rigoroso para testar os limites do Garfo de Hill elástico e a capacidade do MHA de suavizar as oscilações de alta frequência espacial sem perda de consistência energética [47, 48].

A massa linear $M\textit{{{"lin"}}$ (expressa em kg por metro de comprimento axial do cilindro) é calculada diretamente a partir da área da seção transversal da parede do vaso $A}{{"wall"}} = \pi (R\textit{{"ext"}^2 - R}{{"int"}}^2)$ e da densidade média do material [562]. Para os raios de projeto ($R\textit{{{"int"}} = 9,0 "m"$ e $R}{{"ext"}} = 11,0 "m"$), a área do costado resulta em $A\textit{{{"wall"}} approx 125,66 "m"^2$ [562]. A densidade equivalente macroscópica ($hat(\rho)$) do compósito com fração volumétrica simétrica ($V}1 = V_2 = 0,5$) é regida estritamente pela regra aditiva simples [47, 562]:

\begin{equation}
  hat(\rho) = V\textit{1 \rho}{{"aço"}} + V\textit{2 \rho}{{"alu"}} = 0,5 (7800) + 0,5 (2700) = 5250 "kg/m"^3
  \label{eq:densidade_efetiva}
\end{equation}

Resultando em uma massa linear de $659.715 "kg/m"$ para o cilindro bimetálico, o que representa um alívio de peso de exatamente $32,69%$ em comparação a uma estrutura fabricada exclusivamente em aço maciço ($980.148 "kg/m"$), confirmando a elevada eficiência de peso da hibridização de materiais [47, 48].

\subsection{Cenários de Carregamento e Pressão de Projeto}

Para submeter o modelo numérico axissimétrico e a formulação analítica do MHA a condições operacionais realistas e severas, estabeleceu-se uma \textbf{pressão interna de projeto estável de $30 "MPa"$} ($3,0 \times 10^7 "Pa"$) aplicada de forma contínua na fronteira radial interna do cilindro ($r = R_{{"int"}}$) [47, 48, 562].

Esta escolha de carregamento é amplamente fundamentada em aplicações de vasos de pressão modernos na engenharia automotiva e aeroespacial [47]. Conforme ilustrado na literatura técnica e normativa, a faixa de $30$ a $35 "MPa"$ representa a pressão nominal de operação para reservatórios de alta performance Tipo III (que utilizam liner metálico de alumínio ou aço recobertos por enrolamento filamentar composto, CFRP) para o armazenamento seguro de hidrogênio gasoso comprimido ($H_2$) [47, 48].

A imposição deste patamar de solicitação induz um estado de tensões internas elevadas na parede do cilindro de parede grossa ($R_{{"int"}}/t = 4,5$), permitindo-nos avaliar:
1. \textbf{O escoamento inicial das fases:} Onde o aço (com escoamento de $400,0 "MPa"$) e o alumínio (escoamento de $270,0 "MPa"$) são testados sob as tensões equivalentes locais de Von Mises calculadas pelo MEF [47];
2. \textbf{A precisão microscópica do MHA:} Provando se a solução analítica homogeneizada da modelagem assintótica consegue prever as tensões médias e as deformações globais em regime elástico linear mesmo sob pressões internas severas de $300 "bar"$ ($30 "MPa"$) [47, 48].

As condições de contorno de Neumann de pressão interna foram aplicadas no contorno radial interno $Gamma_{{"int"}}$ por meio da forma linear de trabalho virtual descrita na Seção 3.4, garantindo a perfeita integração espacial ponderada pelo raio $r$ exigida pelo FEniCS [48, 562].
